% Intermediate summary matrix for Q2. May later be condensed, moved to an appendix, or redrawn as a figure.
\begin{table}[tbp]
\centering
\footnotesize
\singlespacing
\caption{Working summary for Q2: innate bias versus experience. Working model: intrinsic acoustic bias $\rightarrow$ developmental and associative calibration $\rightarrow$ ecological experience.}
\label{tab:summary-q2}
\begin{tabular}{@{}p{1.3in}p{1.45in}p{2.55in}p{0.85in}@{}}
\toprule
\textbf{System} & \textbf{Na\"ive / early} & \textbf{With experience} & \textbf{Status} \\
\midrule
Avian predator calls & Conflict: defensive in one inbred mouse strain; null in mice, rats, gerbils & Field-experienced rodents respond, form varies by species & Unresolved \\
\addlinespace[2pt]
Rat 22-kHz calls & No selective freezing; frequency-band and neural discrimination present & Selective freezing after aversive experience; self-vocalization not required & Strong \\
\addlinespace[2pt]
Development (sciurids) & Heart-rate discrimination precedes behavior & Call-type physiology maps onto predator class & Supporting \\
\addlinespace[2pt]
Pup calls (mice) & Bias to species-typical intervals & Maternal experience broadens cortical and behavioral tuning & Strong \\
\addlinespace[2pt]
Heterospecific alarms & --- & Eavesdropping widespread; familiarity and acoustic similarity both contribute; safety cues also used & Strong \\
\addlinespace[2pt]
Multipredator / relaxed selection & --- & Framing only; no rodent acoustic population contrast found & Open \\
\bottomrule
\end{tabular}
\end{table}
